\documentclass{article}
\usepackage{graphicx} % Required for inserting images
\usepackage{minted}
\usepackage[a4paper, margin=2cm]{geometry}
\usepackage{tcolorbox}
\usepackage{hyperref}

\title{Tutorato Programmazione 2}
\author{Michele Porcellato, Giacomo Vettore}
\date{8 Aprile 2026}

\begin{document}

\maketitle

\section{Note}
In caso di dubbi, consultare la \href{https://docs.oracle.com/en/java/javase/25/docs/api/index.html}{\textit{documentazione Java}}. Sarà accessibile anche durante l'esame.

\section{Esercizi}
\subsection{SoloVentilatori}
Si progetti \textbf{Solo Ventilatori}, una applicazione dove ventilatori diversi possono pubblicizzarsi, e dove gli utenti possono selezionare a quali ventilatori abbonarsi.
I ventilatori contengono queste informazioni: marca e costo. Inoltre si differenziano per tipologia e alimentazione, e questo influisce sul loro costo di abbonamento, che parte da 10 e si influenza con quanto scritto tra parentesi di sotto.
\newline
Le tipologie di ventilatori sono: 
\begin{itemize}
    \item A piantana (\(+10\))
    \item A parete (\(+20\))
    \item A soffitto (\(+2\))
\end{itemize}
I tipi di alimentazione sono: 
\begin{itemize}
    \item Meccanici (\(\times1\))
    \item Elettrici, che si dividno in 
    \begin{itemize}
        \item A Presa (\(\times2\))
        \item A batteria (\(\times3\))
    \end{itemize}
\end{itemize}
I ventilatori a soffitto \textbf{non} possono essere meccanici.
\newline
In alto viene mostrato il capitale dell'utente (inizializzato a 1000), il costo totale mensile degli abbonamenti (inizializzato a 0) e il numero di mesi trascorsi (inizializzato a 0).
\newline
Dopo di che verranno stampati a video i vari ventilatori.
Per selezionare un ventilatore e abbonarsi, basterà digitare il numero corrispondente.
Per passare al mese successivo, basta cliccare \texttt{Invio}.
Passando al mese successivo, vengono scalati i fondi degli abbonamenti dal capitale dell'utente.
\newline
I tipi di ventilatori saranno i seguenti:
\begin{itemize}
    \item \textbf{Ariete} \\
    Tipo: Piantana Alimentazione: Batteria \\
    Costo: 60
    \item \textbf{Bosch} \\
    Tipo: Parete Alimentazione: Presa \\
    Costo: 60
    \item \textbf{Bosch} \\
    Tipo: Parete Alimentazione: Meccanica \\
    Costo: 30
    \item \textbf{Bosch} \\
    Tipo: Soffitto Alimentazione: Batteria \\
    Costo: 36
    \item \textbf{Parkside} \\
    Tipo: Soffitto Alimentazione: Presa \\
    Costo: 24
\end{itemize}
Se l'utente non ha abbastanza fondi, viene stampato un messaggio di errore sulla console.
\begin{itemize}
    \item Si usi, quando evidente/utile, una gerarchia di ereditarietà.
    \item Si usino costanti ove ragionevole, e si badi alla pulizia del codice, evitando duplicazioni e codice inutilmente complesso.
\end{itemize}

\subsection*{Esempio 1}
\begin{tcolorbox}
---------------------------------------- \\
Capitale: 1000 \\
Costo totale mensile abbonamenti: 0 \\
Mesi trascorsi: 0 \\
---------------------------------------- \\
Ventilatori disponibili: \\
1) Ariete \\
Tipo: Piantana Alimentazione: Batteria \\
Costo: 60 \\
2) Bosch \\
Tipo: Parete Alimentazione: Presa \\
Costo: 60 \\
3) Bosch \\
Tipo: Parete Alimentazione: Meccanica \\
Costo: 30 \\
4) Bosch \\
Tipo: Soffitto Alimentazione: Batteria \\
Costo: 36 \\
5) Parkside \\
Tipo: Soffitto Alimentazione: Presa \\
Costo: 24 \\
---------------------------------------- \\
Seleziona il numero di un ventilatore per abbonarti. \\
Premi Invio senza scrivere nulla per passare al mese successivo. \\
\(>\) \texttt{1} \\
Abbonamento aggiunto. Nuovo costo mensile: 60. \\
\end{tcolorbox}

\subsection*{Esempio 2}
\begin{tcolorbox}
---------------------------------------- \\
Capitale: 160 \\
Costo totale mensile abbonamenti: 210 \\
Mesi trascorsi: 4 \\
---------------------------------------- \\
Ventilatori disponibili: \\
1) Ariete \\
Tipo: Piantana Alimentazione: Batteria \\
Costo: 60 \\
2) Bosch \\
Tipo: Parete Alimentazione: Presa \\
Costo: 60 \\
3) Bosch \\
Tipo: Parete Alimentazione: Meccanica \\
Costo: 30 \\
4) Bosch \\
Tipo: Soffitto Alimentazione: Batteria \\ 
Costo: 36 \\
5) Parkside \\
Tipo: Soffitto Alimentazione: Presa \\
Costo: 24 \\
---------------------------------------- \\
Seleziona il numero di un ventilatore per abbonarti. \\
Premi Invio senza scrivere nulla per passare al mese successivo. \\
\(>\) \\
Errore: fondi insufficienti per pagare gli abbonamenti mensili.\\
\end{tcolorbox}

\subsection*{Suggerimenti}
La soluzione proposta usa MVC. All'esame sarà richiesto ma per ora è possibile svolgere l'esercizio anche senza.
\newline
Si consiglia di implementare una classe \texttt{Utente}, per gestire il saldo e una classe \texttt{Vetrina}, per mostrare i ventilatori disponibili.

\end{document}
